% GENERATED FILE. Edit canonical lessons, then run npm run generate:books.
% canonical-source-hash: fnv1a64:19c152e0a1b620ce
% canonical-lessons: GE-C496-stillen, GE-C496-windel, GE-C496-wickeln, GE-C496-kinderwagen, GE-C496-schnuller

\chapter{Ein Baby im Haus}
\label{ch:ge-ein-baby-im-haus}
\hlchaptermodality{496}

\begin{chapteropening}
\textbf{By the end of this chapter:} \emph{I can talk about looking after a baby.}

Complete the last lesson of chapter 496: I can talk about looking after a baby.
\end{chapteropening}

\section[stillen]{stillen — to breastfeed; to quench}

\label{lesson:GE-C496-stillen}



\begin{quote}
Before the new one: say the German for to obey, then the German for puberty.
\end{quote}

\subsection*{You'll want to know: stillen}
\textbf{stillen} — \textquotedblleft{}to breastfeed; to quench\textquotedblright{}.

\begin{quote}
Sie stillt ihr Baby alle drei Stunden.
\end{quote}

\begin{quote}
\emph{She breastfeeds her baby every three hours.}
\end{quote}

Also to quench: \emph{den Durst stillen}.

\begin{cousinweb}
\emph{stillen} is built on \emph{still} \textquotedblleft{}quiet, calm\textquotedblright{} (EtymDB): first to make quiet, then to calm a baby at the breast.
\end{cousinweb}

\begin{grammarlens}[title={What you've built}]
One more word for telling, explaining and describing at B1.
\end{grammarlens}

\subsection*{Your turn}
\begin{itemize}
\raggedright
  \item \emph{Say it:} \emph{stillen}
  \item \emph{Say it:} \emph{stillen}, once more
  \item \emph{Say it:} a sentence of your own with \emph{stillen}
  \item \emph{Recall:} say the German for to obey, then the German for puberty, then say \emph{stillen} again
\end{itemize}

\begin{tcolorbox}[breakable,colback=teal!4,colframe=teal!35!black,title={Before you move on}]
What does stillen mean? (\textquotedblleft{}To breastfeed; to quench\textquotedblright{}.) Say it once more.
\end{tcolorbox}
\section[die Windel]{die Windel — nappy, diaper}

\label{lesson:GE-C496-windel}



\begin{quote}
Before the new one: say the German for puberty, then the German for to breastfeed.
\end{quote}

\subsection*{You'll want to know: die Windel}
\textbf{die Windel} — \textquotedblleft{}nappy, diaper\textquotedblright{}.

\begin{quote}
Das Baby braucht eine frische Windel.
\end{quote}

\begin{quote}
\emph{The baby needs a fresh nappy.}
\end{quote}

Plural: \emph{die Windeln}.

\begin{cousinweb}
\emph{Windel} is inherited from Middle High German \emph{windel} and Old High German \emph{windila} (EtymDB), built on \emph{winden} \textquotedblleft{}to wind, to wrap\textquotedblright{}.
\end{cousinweb}

\begin{grammarlens}[title={What you've built}]
One more word for telling, explaining and describing at B1.
\end{grammarlens}

\subsection*{Your turn}
\begin{itemize}
\raggedright
  \item \emph{Say it:} \emph{die Windel}
  \item \emph{Say it:} \emph{die Windel}, once more
  \item \emph{Say it:} a sentence of your own with \emph{die Windel}
  \item \emph{Recall:} say the German for puberty, then the German for to breastfeed, then say \emph{die Windel} again
\end{itemize}

\begin{tcolorbox}[breakable,colback=teal!4,colframe=teal!35!black,title={Before you move on}]
What does die Windel mean? (\textquotedblleft{}Nappy, diaper\textquotedblright{}.) Say it once more.
\end{tcolorbox}
\section[wickeln]{wickeln — to change (a baby); to wrap}

\label{lesson:GE-C496-wickeln}



\begin{quote}
Before the new one: say the German for to breastfeed, then the German for nappy.
\end{quote}

\subsection*{You'll want to know: wickeln}
\textbf{wickeln} — \textquotedblleft{}to change (a baby); to wrap\textquotedblright{}.

\begin{quote}
Der Vater wickelt das Baby im Bad.
\end{quote}

\begin{quote}
\emph{The father changes the baby in the bathroom.}
\end{quote}

Also to wrap: \emph{ein Geschenk in Papier wickeln}.

\begin{cousinweb}
\emph{wickeln} is inherited from Middle High German \emph{wickeln} \textquotedblleft{}to wrap around\textquotedblright{} (EtymDB).
\end{cousinweb}

\begin{grammarlens}[title={What you've built}]
One more word for telling, explaining and describing at B1.
\end{grammarlens}

\subsection*{Your turn}
\begin{itemize}
\raggedright
  \item \emph{Say it:} \emph{wickeln}
  \item \emph{Say it:} \emph{wickeln}, once more
  \item \emph{Say it:} a sentence of your own with \emph{wickeln}
  \item \emph{Recall:} say the German for to breastfeed, then the German for nappy, then say \emph{wickeln} again
\end{itemize}

\begin{tcolorbox}[breakable,colback=teal!4,colframe=teal!35!black,title={Before you move on}]
What does wickeln mean? (\textquotedblleft{}To change (a baby); to wrap\textquotedblright{}.) Say it once more.
\end{tcolorbox}
\section[der Kinderwagen]{der Kinderwagen — pram, baby buggy}

\label{lesson:GE-C496-kinderwagen}



\begin{quote}
Before the new one: say the German for nappy, then the German for to change.
\end{quote}

\subsection*{You'll want to know: der Kinderwagen}
\textbf{der Kinderwagen} — \textquotedblleft{}pram, baby buggy\textquotedblright{}.

\begin{quote}
Mit dem Kinderwagen gehen wir jeden Tag in den Park.
\end{quote}

\begin{quote}
\emph{We go to the park with the pram every day.}
\end{quote}

Plural: \emph{die Kinderwagen}.

\begin{cousinweb}
\emph{Kinderwagen} is built on \emph{Kind} \textquotedblleft{}child\textquotedblright{} and \emph{Wagen} \textquotedblleft{}cart\textquotedblright{}, inherited from Old High German \emph{wagan}, from Proto-Germanic \emph{wagnaz} (EtymDB).
\end{cousinweb}

\begin{grammarlens}[title={What you've built}]
One more word for telling, explaining and describing at B1.
\end{grammarlens}

\subsection*{Your turn}
\begin{itemize}
\raggedright
  \item \emph{Say it:} \emph{der Kinderwagen}
  \item \emph{Say it:} \emph{der Kinderwagen}, once more
  \item \emph{Say it:} a sentence of your own with \emph{der Kinderwagen}
  \item \emph{Recall:} say the German for nappy, then the German for to change, then say \emph{der Kinderwagen} again
\end{itemize}

\begin{tcolorbox}[breakable,colback=teal!4,colframe=teal!35!black,title={Before you move on}]
What does der Kinderwagen mean? (\textquotedblleft{}Pram, baby buggy\textquotedblright{}.) Say it once more.
\end{tcolorbox}
\section[der Schnuller]{der Schnuller — dummy, pacifier}

\label{lesson:GE-C496-schnuller}



\begin{quote}
Before the new one: say the German for to change, then the German for pram.
\end{quote}

\subsection*{You'll want to know: der Schnuller}
\textbf{der Schnuller} — \textquotedblleft{}dummy, pacifier\textquotedblright{}.

\begin{quote}
Ohne Schnuller schläft das Baby nicht ein.
\end{quote}

\begin{quote}
\emph{The baby won't fall asleep without a dummy.}
\end{quote}

Plural: \emph{die Schnuller}. In Switzerland: \emph{der Nuggi}.

\begin{cousinweb}
\emph{Schnuller} is usually traced to the dialect verb \emph{schnullen} \textquotedblleft{}to suck\textquotedblright{}; EtymDB has no entry for it.
\end{cousinweb}

\begin{grammarlens}[title={What you've built}]
One more word for telling, explaining and describing at B1.
\end{grammarlens}

\subsection*{Your turn}
\begin{itemize}
\raggedright
  \item \emph{Say it:} \emph{der Schnuller}
  \item \emph{Say it:} \emph{der Schnuller}, once more
  \item \emph{Say it:} a sentence of your own with \emph{der Schnuller}
  \item \emph{Recall:} say the German for to change, then the German for pram, then say \emph{der Schnuller} again
\end{itemize}

\begin{tcolorbox}[breakable,colback=teal!4,colframe=teal!35!black,title={Before you move on}]
What does der Schnuller mean? (\textquotedblleft{}Dummy, pacifier\textquotedblright{}.) Say it once more.
\end{tcolorbox}
